\chapter{Lists, Quotations, and Notes}
\label{chap:lists}
\index{list!bulleted}\index{list!numbered}\index{list!description}\index{quotation}

Methods often contain criteria; reports contain recommendations; teaching
documents contain steps. A list can make those relationships visible, but a
list used for every pair of sentences fragments the argument. \LaTeX{} provides
structures for unordered, ordered, and labelled items, plus distinct forms for
short and multi-paragraph quotations.

\begin{learningoutcomes}
  \item Create unordered, ordered, and description lists.
  \item Nest lists without losing their logical hierarchy.
  \item Choose between the \code{quote} and \code{quotation} environments.
  \item Insert footnotes at suitable locations.
  \item Explain when margin notes are unsafe or unnecessary.
  \item Customise list spacing conservatively with \pkg{enumitem}.
\end{learningoutcomes}

\section{Unordered lists with itemize}

Use \code{itemize} when sequence does not change meaning:

\begin{latexcode}{Complete example: an unordered list}
\documentclass{article}
\begin{document}
The comparison requires:
\begin{itemize}
  \item paired measurements;
  \item a recorded unit; and
  \item an explicit missing-data rule.
\end{itemize}
\end{document}
\end{latexcode}

Each \cmd{item} begins an item. The punctuation is part of the sentence
surrounding the list: these fragments use semicolons and the last ends with a
full stop. Complete-sentence items can each end with a full stop instead.

\section{Ordered lists with enumerate}

Use \code{enumerate} when order, rank, or later reference matters:

\begin{latexcode}{Complete example: a procedure}
\begin{enumerate}
  \item Record the Method A value.
  \item Record the Method B value.
  \item Calculate the paired difference.
  \item Inspect the result before summarising it.
\end{enumerate}
\end{latexcode}

This fragment belongs inside a document body. \LaTeX{} supplies the numbers and
renumbers later items when one is inserted. If the prose must refer repeatedly
to a particular item, attach a label rather than typing its number; labels are
introduced in Chapter~\ref{chap:references}.

\begin{pausepredict}
Add a new second item. Predict the number of the original second item and every
item after it. What maintenance problem would arise if the numbers were typed
as ordinary text?
\end{pausepredict}

\section{Description lists}

A description list pairs a term with its explanation:

\begin{latexcode}{Complete example: labelled definitions}
\begin{description}
  \item[Method A] The reference teaching method.
  \item[Method B] The comparison teaching method.
  \item[Difference] Method B minus Method A for one pair.
\end{description}
\end{latexcode}

The text in square brackets is the item label, not an automatically numbered
optional setting. Keep labels short enough to remain readable. For a large
alphabetical vocabulary, a glossary is more suitable than a very long list in
the middle of a chapter.

\section{Nested lists}

A list can contain another list when the hierarchy is genuine:

\begin{latexcode}{Complete example: one nested level}
\begin{enumerate}
  \item Validate the input.
    \begin{itemize}
      \item Check that each pair has an identifier.
      \item Check that both values use the same unit.
    \end{itemize}
  \item Calculate the differences.
\end{enumerate}
\end{latexcode}

Indentation helps readers of the source see which environment is inside which.
Deeply nested lists are difficult to read in source and PDF. If the hierarchy
needs three or four levels, consider headings, a table, or separate procedural
stages.

\begin{commonerror}
\textbf{Symptom:} ``Something's wrong--perhaps a missing \cmd{item}''.

\textbf{Likely cause:} text has been placed directly inside a list before its
first \cmd{item}, or an item command is missing.

\textbf{Repair:} inspect the first lines after \cmd{begin}\required{itemize} or
\cmd{begin}\required{enumerate}. Make every list entry begin with \cmd{item}.
\end{commonerror}

\section{Short and long quotations}

Use the \code{quote} environment for a short displayed quotation. Use
\code{quotation} when the quoted material contains multiple paragraphs; it
marks their paragraph structure.

\begin{latexcode}{Displayed quotation pattern}
The protocol defines repeatability as follows:
\begin{quote}
  [Insert a short, verified quotation and cite its source.]
\end{quote}
The analysis applies that definition to simulated values.
\end{latexcode}

The bracketed sentence is a placeholder, not a quotation. Do not invent words
and place them in quotation marks. When exact wording is needed, check the
source, quote only what is necessary, preserve its meaning, and cite it. Most
scholarly explanation should paraphrase responsibly in the author's own voice.

\begin{scienceinpractice}
A standards report may need one exact definition because later requirements
depend on its wording. A literature review rarely needs long blocks of copied
text. In both cases, quotation formatting is separate from permission,
accuracy, and citation responsibility.
\end{scienceinpractice}

\section{Footnotes and optional margin notes}

Footnotes are useful for a concise qualification that would interrupt the
main sentence. Put the command after the word or punctuation associated with
the note, following the chosen style consistently. Do not put a full data table
or indispensable method in a footnote.

Margin notes require a package or class command and can collide with floats,
two-column layouts, or binding margins. They are therefore optional advanced
material, not a default scholarly annotation system.

\begin{advancedtopic}
Some classes provide \cmd{marginpar}\required{text}. Before using it, check
whether the document is one- or two-sided and whether the target venue permits
margin material. Draft reminders are often safer as source comments because
they cannot accidentally become publication content.
\end{advancedtopic}

\section{Careful customisation with enumitem}

The \pkg{enumitem} package can adjust labels and spacing. Load it in the
preamble, then state a local purpose:

\begin{latexcode}{Complete example: a compact local list}
\documentclass{article}
\usepackage{enumitem}
\begin{document}
\begin{enumerate}[label=Step \arabic*:,itemsep=0.25em]
  \item Import the paired values.
  \item Confirm the unit.
  \item Calculate differences.
\end{enumerate}
\end{document}
\end{latexcode}

The option changes this list only. Avoid setting every list to zero spacing;
the apparent saving can make a dense technical page harder to scan. A class
may already define list spacing appropriate to its layout.

\begin{goodpractice}
Choose a list because the relationship is unordered, sequential, or
term-and-definition. Keep parallel grammatical form: do not mix ``Record the
value'', ``Calculation of the mean'', and ``The plot is inspected'' in one
procedural list.
\end{goodpractice}

\chapterproject

\begin{miniproject}
In the Methods section of the running article, add a three-item ordered list
describing data validation, difference calculation, and inspection. Add an
unordered list of two inclusion criteria. Keep the items grammatically
parallel. Add a source comment for a future verified definition; do not insert
an invented quotation.
\end{miniproject}

\chaptersummary

\code{itemize} expresses an unordered set, \code{enumerate} a sequence, and
\code{description} a labelled relationship. Nested lists should reflect a
real, shallow hierarchy. \code{quote} suits a short displayed quotation;
\code{quotation} supports multiple quoted paragraphs. Footnotes hold concise
supplements, while margin notes are optional and layout-sensitive.
\pkg{enumitem} can customise a list when the reason is explicit.

\chaptercheck
\begin{chapterchecklist}
  \item I chose the list type from the relationship between items.
  \item Every entry begins with \cmd{item}.
  \item Nested environments close in the reverse order they opened.
  \item My list items use parallel grammar and consistent punctuation.
  \item Any exact quotation has a verified source.
  \item Important methods and evidence remain in the main text.
\end{chapterchecklist}

\chapterexercisestitle
\begin{chapterexercises}
  \item Choose a list type for materials, a protocol, and a glossary of three
        terms.
  \item Convert three manually numbered lines into \code{enumerate}.
  \item Add one nested checklist under the first procedural step.
  \item Repair a list whose first entry lacks \cmd{item}.
  \item Explain the difference between \code{quote} and \code{quotation}.
  \item Rewrite a list so that all items begin with active verbs.
  \item Give a reason to keep a draft note as a source comment rather than a
        margin note.
  \item Challenge: redesign a dense list from a report as either prose, a
        table, or two smaller lists, and justify the choice.
\end{chapterexercises}
